% Generated draft: requires mathematical and source review.
\documentclass[11pt]{article}
\usepackage[margin=25mm]{geometry}
\usepackage{fontspec}
\setmainfont{Latin Modern Roman}
\usepackage{amsmath,amssymb,mathtools,unicode-math}
\setmathfont{Latin Modern Math}
\usepackage{enumitem,needspace,xurl,hyperref}
\setlength{\parindent}{0pt}
\setlength{\parskip}{5pt}
\setlength{\emergencystretch}{3em}
\newcommand{\N}{\mathbb{N}}
\newcommand{\Z}{\mathbb{Z}}
\newcommand{\Q}{\mathbb{Q}}
\newcommand{\R}{\mathbb{R}}
\newcommand{\C}{\mathbb{C}}
\newcommand{\F}{\mathbb{F}}
\newcommand{\A}{\mathbb{A}}
\newcommand{\PP}{\mathbb{P}}
\newcommand{\E}{\mathbb{E}}
\newcommand{\eps}{\varepsilon}
\newcommand{\dd}{\,\mathrm{d}}
\begin{document}
\section*{UnsolvedMath ID: 20002293}
\textbf{Source problem code:} catalogue entry
\subsection{Definitions and mathematical statement}
For finite \(n\), put \(K_n=\Q(t_1,\ldots,t_n)\), with the variables algebraically independent. Flenner's Question \(9\) asks for one sentence \(\varphi\) in the pure ring language such that
\[
 K_n\models\varphi\quad\Longleftrightarrow\quad n\text{ is even}.
\]
The source does not state whether \(n=0\), \(K_0=\Q\), is included; either convention must be used consistently for the even/odd families below. It suggests ultrapowers as a possible negative approach, not as a completed proof. The aid reports the parity question open.

The aid's specialization statement is that every finite \(K_n\) has the same existential and universal ring theories as \(\Q\). A finite existential witness uses finitely many rational functions; specializing the variables to rational values outside the finitely many denominator and inequation zero sets preserves its quantifier-free requirements. Conversely rational witnesses remain witnesses in \(K_n\). Universal agreement follows by negating existential sentences. Thus no existential or universal sentence can define the parity, but arbitrary alternating-quantifier sentences remain outside this argument.

For the full question, fix an enumeration \(\sigma_1,\sigma_2,\ldots\) of ring sentences. The aid's exact finite-fragment criterion is
\[
 \text{no parity-defining sentence exists}
 \quad\Longleftrightarrow\quad
 \forall N\ \exists e\text{ even},\ o\text{ odd}:\quad
 K_e,K_o\text{ agree on }\sigma_1,\ldots,\sigma_N.
\]
Equivalently they agree on arbitrarily large finite sentence fragments, or on each finite sentence-length bound under a fixed finite encoding. If some finite fragment had disjoint sets of even and odd truth types, the finite disjunction of the even types, using sentences or their negations, would itself give a parity sentence.

If these pairs \((e_N,o_N)\) exist, for any nonprincipal ultrafilter \(\mathcal U\) on \(\omega\), Lo\'s's theorem yields
\[
 \prod_N K_{e_N}/\mathcal U\ \equiv\
 \prod_N K_{o_N}/\mathcal U.
\]
Conversely such elementary equivalence of an all-even and an all-odd varying-field ultraproduct would rule out a parity sentence. These are ultraproducts of possibly varying fields, not a claim that taking an ultrapower of one fixed field changes its theory. The reduction does not prove existence of the required finite-fragment pairs or of the displayed equivalent ultraproducts, and therefore does not solve the full parity problem.
\subsection{Short English statement}
All finite rational function fields share the existential and universal theories of \(\Q\), so those sentence classes cannot detect parity. The aid reduces full undefinability to even/odd agreement on every finite fragment, equivalently a suitable pair of elementarily equivalent varying-field ultraproducts. It does not establish that agreement.
\subsection{Sources}
\begin{itemize}
\item Local UnsolvedMath registry, permanent ID \(20002293\), source code \texttt{catalogue entry}; canonical collection \texttt{aim-logic-notes.json}, notes[68].
\item AIM workshop ``Definability and decidability problems in number theory,'' Question \(9\), Flenner: \url{https://aimath.org/pastworkshops/definabilityinntproblems.pdf}. The supplied verification reports restoration of subscripts and punctuation from the source mathematical fonts. No PDF was opened in this session.
\item The aid uses specialization and Lo\'s's ultraproduct theorem without full authoritative bibliographic entries. Original remarks and literature fields are empty; no proposer-field value or year is supplied.
\item Supplied research attempt \(1\), result ``reduction,'' run status ``open,'' rigor label ``proved,'' novelty confidence ``low.'' Registry status \texttt{partially\_solved}, classification \texttt{PARTIAL-PROGRESS}, suggested difficulty \texttt{L4}, updated 14 August 2026.
\end{itemize}
\end{document}
